
\chapter{如期完成全部论文工作的可能性}
本课题目前已成功完成多模态蛋白质数据的收集与高质量预处理，涵盖蛋白质序列、结构与功能标签等多源异构信息。在此基础上，构建了融合蛋白质序列语义、空间结构的异构图表示模型，并基于异构图神经网络实现了统一且高效的多模态预训练表示学习架构。在建模与实验方面，课题已搭建完整的多任务训练框架。

初步实验结果表明，该模型在分子功能（MFO）、细胞成分（CCO）等GO子本体任务展现出较强的预测能力与语义泛化能力，尤其在零样本功能预测场景中取得了有竞争力的性能。目前已有初步成果，正在整理实验结果，计划近期撰写英文论文投稿至人工智能、生物信息学方向的国际会议或高水平期刊。后续将进一步完善语义匹配机制，分别优化编码器与分类器性能，引入GO标签的层次结构辅助监督，并拓展在长尾标签与跨物种数据上的评估。

综上所述，课题整体技术路线清晰、实施方案可行，阶段性成果扎实，为后续工作的高效推进奠定了坚实基础。预计在既定时间范围内可如期完成全部研究任务，并有望产出具有理论价值与实际应用前景的创新性研究成果。


